\documentclass[../../main.tex]{subfiles}
\begin{document}
\section{余向量}\label{section:ch11-余向量}

设 $V$ 为有限维向量空间。（照例，我们的向量空间都假定是实的。）我们把 $V$ 上的\textbf{余向量}定义为 $V$ 上的实值线性泛函，即线性映射 $\omega\colon V\to\mathbb{R}$。$V$ 上所有余向量的空间，在逐点加法与数乘的显然运算下，本身是一个实向量空间。它记为 $V^*$，称为 $V$ 的\textbf{对偶空间}。

下面的命题给出了有限维情形中关于 $V^*$ 的最重要事实。回顾习题~\ref{exercise:appB-13}：线性映射由它在任意一组基上的值唯一确定。

\begin{proposition}\label{proposition:ch11-对偶基}
设 $V$ 为有限维向量空间。给定 $V$ 的任意一组基 $(E_1,\dots,E_n)$，令 $\varepsilon^1,\dots,\varepsilon^n\in V^*$ 为由
\begin{align*}
\varepsilon^i(E_j)=\delta^i_j
\end{align*}
定义的余向量，其中 $\delta^i_j$ 是由\eqref{align:ch04-kronecker}~定义的 Kronecker delta 符号。则 $(\varepsilon^1,\dots,\varepsilon^n)$ 是 $V^*$ 的一组基，称为 $(E_j)$ 的\textbf{对偶基}。因此 $\dim V^*=\dim V$。
\end{proposition}

\begin{exercise}\label{exercise:ch11-02}
证明命题~\ref{proposition:ch11-对偶基}。
\end{exercise}

例如，把此命题应用于 $\mathbb{R}^n$ 的标准基 $(e_1,\dots,e_n)$。其对偶基记为 $(e^1,\dots,e^n)$（注意上指标），称为\textbf{标准对偶基}。这些基余向量是 $\mathbb{R}^n$ 上由
\begin{align*}
e^i(v)=e^i(v^1,\dots,v^n)=v^i
\end{align*}
给出的线性泛函。换言之，$e^i$ 是取出向量第 $i$ 个分量的线性泛函。在矩阵记号下，从 $\mathbb{R}^n$ 到 $\mathbb{R}$ 的线性映射由 $1\times n$ 矩阵表示，称为行矩阵。因此基余向量也可视为由行矩阵
\begin{align*}
e^1=(1\ 0\ \cdots\ 0),\quad e^2=(0\ 1\ 0\ \cdots\ 0),\quad \dots,\quad e^n=(0\ \cdots\ 0\ 1)
\end{align*}
表示的线性泛函。

一般地，若 $(E_j)$ 是 $V$ 的一组基，$(\varepsilon^i)$ 是其对偶基，则对任意向量 $v=v^jE_j\in V$（使用 Einstein 求和约定），有
\begin{align*}
\varepsilon^i(v)=v^j\varepsilon^i(E_j)=v^j\delta^i_j=v^i.
\end{align*}
于是，正如在 $\mathbb{R}^n$ 的情形中一样，第 $i$ 个基余向量 $\varepsilon^i$ 取出向量 $v$ 的第 $i$ 个分量。

由对偶基 $(\varepsilon^i)$，任意余向量 $\omega\in V^*$ 可写为
\begin{align*}
\omega=\omega_i\varepsilon^i,
\end{align*}
其中系数 $(\omega_1,\dots,\omega_n)$ 由
\begin{align*}
\omega_i=\omega(E_i)
\end{align*}
给出。若 $v=v^jE_j$，则 $\omega(v)=\omega_i v^i$。

我们经常把余向量 $\omega\in V^*$ 作用于向量 $v\in V$ 的结果写成一个“括号” $\langle\omega,v\rangle$，并称之为 $\omega$ 与 $v$ 的\textbf{自然配对}：
\begin{align}
\langle\omega,v\rangle=\omega(v).\label{align:ch11-自然配对}
\end{align}

余向量与流形的相关性如下。设 $M$ 为光滑流形。对每一点 $p\in M$，$T_pM$ 的对偶空间 $T_p^*M$ 称为 $p$ 处的\textbf{余切空间}。其元素（即 $T_pM$ 上的线性泛函）称为 $p$ 处的\textbf{切余向量}或\textbf{余切向量}。所有余切空间的不交并
\begin{align*}
T^*M=\coprod_{p\in M}T_p^*M
\end{align*}
称为 $M$ 的\textbf{余切丛}。

我们暂时在这一章的前几节中，把余向量场定义为余切丛的截面：即映射 $\eta\colon M\to T^*M$，使得对每个 $p\in M$ 有 $\eta(p)\in T_p^*M$。我们将看到，这正是我们所需要的几何对象。在第~12~章中，我们将把这一概念与流形上张量场的一般定义统一起来。

设 $(U,\varphi)$ 是 $M$ 的光滑坐标图，坐标为 $(x^i)$。对每个 $p\in U$，切空间 $T_pM$ 有坐标基 $(\partial/\partial x^i|_p)$。我们记相应的对偶基为 $(\theta^1|_p,\dots,\theta^n|_p)$，其特征是
\begin{align}
\theta^i|_p\left(\frac{\partial}{\partial x^j}\bigg|_p\right)=\delta^i_j.\label{align:ch11-坐标余基}
\end{align}
把 $p$ 变动，就得到 $U$ 上 $T^*M$ 的一组 $n$ 个截面 $(\theta^1,\dots,\theta^n)$。

\begin{proposition}\label{proposition:ch11-坐标余基标架}
设 $M$ 为带边或不带边的光滑流形。对 $M$ 的每个光滑坐标图 $(U,\varphi)$，相应的坐标余向量场 $(\theta^1,\dots,\theta^n)$ 是 $U$ 上余切丛的一个局部标架。
\end{proposition}

\begin{proof}
由构造，对每个 $p\in U$，$(\theta^1|_p,\dots,\theta^n|_p)$ 是 $T_p^*M$ 的一组基。
\end{proof}

为了计算坐标余向量场在坐标变换下如何变化，设 $(U,\varphi)$ 与 $(V,\psi)$ 是两个光滑坐标图，坐标分别为 $(x^i)$ 与 $(\tilde x^j)$，且 $U\cap V\neq\varnothing$。由\cref{proposition:ch11-对偶基}，新的余基 $(\tilde\theta^j)$ 可用旧的余基 $(\theta^i)$ 表示。利用\eqref{align:ch11-坐标余基}，我们看到
\begin{align*}
\tilde\theta^j\left(\frac{\partial}{\partial x^i}\bigg|_p\right)=\tilde\theta^j\left(\frac{\partial\tilde x^k}{\partial x^i}(p)\frac{\partial}{\partial\tilde x^k}\bigg|_p\right)=\frac{\partial\tilde x^j}{\partial x^i}(p).
\end{align*}
（这与向量的变换公式相似但转置了。）这表明
\begin{align}
\tilde\theta^j=\frac{\partial\tilde x^j}{\partial x^i}\theta^i.\label{align:ch11-余向量坐标变换}
\end{align}

若 $\eta$ 是 $U$ 上的光滑余向量场，我们可写 $\eta=\eta_i\theta^i$，其中 $\eta_i\colon U\to\mathbb{R}$ 是分量函数
\begin{align*}
\eta_i(p)=\left\langle\eta_p,\frac{\partial}{\partial x^i}\bigg|_p\right\rangle.
\end{align*}
由\eqref{align:ch11-余向量坐标变换}，分量在坐标变换下的变换规则为
\begin{align*}
\tilde\eta_j=\eta_i\frac{\partial x^i}{\partial\tilde x^j}.
\end{align*}

设 $M$ 为光滑流形。$T^*M$ 的一个\textbf{光滑余向量场}是余切丛的光滑截面。我们记 $M$ 上光滑余向量场的空间为 $\mathfrak X^*(M)$。（这是另一个“协变向量”一词虽不幸、却已根深蒂固的例子。）

\begin{example}\label{example:ch11-余向量场例子}
在 $\mathbb{R}^n$ 上，余切丛自然平凡化为 $\mathbb{R}^n\times(\mathbb{R}^n)^*$，而 $(\mathbb{R}^n)^*$ 又自然等同为 $\mathbb{R}^n$。于是余向量场可等同于从 $\mathbb{R}^n$ 到 $\mathbb{R}^n$ 的映射。若坐标为 $(x^1,\dots,x^n)$，则标准余基为 $(e^1,\dots,e^n)$，任意余向量场可写为 $\omega=\omega_i e^i$，其中 $\omega_i$ 是 $x$ 的函数。光滑性判据与向量场完全类似。$\blacktriangleright$
\end{example}

下面的命题给出了光滑余向量场的坐标刻画。

\begin{proposition}[光滑余向量场的坐标判别法]\label{proposition:ch11-光滑余向量场坐标判别}
设 $M$ 为带边或不带边的光滑流形，$\eta$ 为 $M$ 上的余向量场。则 $\eta$ 光滑，当且仅当对 $M$ 的每个光滑坐标图 $(U,(x^i))$，其分量函数 $\eta_i=\langle\eta,\partial/\partial x^i\rangle$ 在 $U$ 上都是光滑的。
\end{proposition}

\begin{proof}
设 $\eta$ 在 $U$ 上的分量函数为 $\eta_i$。若 $\eta$ 是光滑截面，则每个 $\eta_i=\langle\eta,\partial/\partial x^i\rangle$ 是光滑函数与光滑向量场的自然配对，因而是光滑函数。

反过来，设每个坐标图上的分量函数都光滑。要证 $\eta$ 是 $T^*M$ 的光滑截面，只需证明它在每个坐标图 $U$ 上光滑。由定义，$T^*M$ 在 $U$ 上的坐标是由坐标余基 $(\theta^i)$ 决定的向量丛坐标 $(x^i, v_i)$，而 $\eta$ 在此坐标下为
\begin{align*}
\eta(x)=(x^1,\dots,x^n,\eta_1(x),\dots,\eta_n(x)).
\end{align*}
由于分量函数光滑，$\eta$ 在 $U$ 上光滑。
\end{proof}

余向量在光滑映射下的行为与向量场相反。设 $F\colon M\to N$ 是光滑流形间的光滑映射，$p\in M$。微分 $dF_p\colon T_pM\to T_{F(p)}N$ 的对偶是一个线性映射，它把 $F(p)$ 处的余向量映到 $p$ 处的余向量：
\begin{align*}
dF_p^*\colon T_{F(p)}^*N\to T_p^*M.
\end{align*}
若 $\omega\in T_{F(p)}^*N$，其\textbf{拉回} $dF_p^*\omega\in T_p^*M$ 定义为
\begin{align}
dF_p^*\omega(v)=\omega(dF_p(v)),\quad v\in T_pM.\label{align:ch11-余向量拉回}
\end{align}
注意我们在 $dF_p$ 上加星号来表示对偶映射。（它不是 $dF_p$ 的伴随——这里没有内积——只是对偶线性映射。）

与向量场不同，余向量场总是可以被光滑映射拉回。若 $F\colon M\to N$ 是光滑映射，而 $\omega$ 是 $N$ 上的光滑余向量场，我们定义 $M$ 上的余向量场 $F^*\omega$ 为
\begin{align}
(F^*\omega)_p=dF_p^*(\omega_{F(p)}).\label{align:ch11-余向量场拉回}
\end{align}
下面的命题表明它是光滑的。

\begin{proposition}\label{proposition:ch11-拉回余向量场光滑}
设 $F\colon M\to N$ 为光滑流形间的光滑映射。若 $\omega\in\mathfrak X^*(N)$ 是光滑余向量场，则 $F^*\omega\in\mathfrak X^*(M)$ 是光滑的。
\end{proposition}

\begin{proof}
我们在坐标下检验。若 $(x^i)$ 是 $M$ 上的坐标，$(y^j)$ 是 $N$ 上的坐标，且 $\omega=\omega_j\,dy^j$，则
\begin{align*}
(F^*\omega)_p(v)=\omega_{F(p)}(dF_p(v))=\omega_j(F(p))\,dy^j(dF_p(v))=\omega_j(F(p))\frac{\partial F^j}{\partial x^i}(p)v^i.
\end{align*}
因此 $F^*\omega$ 的第 $i$ 个分量为
\begin{align*}
(F^*\omega)_i(p)=\omega_j(F(p))\frac{\partial F^j}{\partial x^i}(p),
\end{align*}
它是光滑函数。由命题~\ref{proposition:ch11-光滑余向量场坐标判别}，$F^*\omega$ 光滑。
\end{proof}

下面的命题总结了拉回的性质。它的证明是直接的，留作习题。

\begin{proposition}\label{proposition:ch11-拉回性质}
设 $F\colon M\to N$ 与 $G\colon N\to P$ 为光滑流形间的光滑映射。则
\begin{enumerate}[(a)]
\item $F^*\colon\mathfrak X^*(N)\to\mathfrak X^*(M)$ 是线性映射；
\item 对任意 $f\in C^\infty(N)$ 与 $\omega\in\mathfrak X^*(N)$，有 $F^*(f\omega)=(f\circ F)F^*\omega$；
\item $(G\circ F)^*\omega=F^*(G^*\omega)$，对任意 $\omega\in\mathfrak X^*(P)$；
\item $\mathrm{Id}^*\omega=\omega$，对任意 $\omega\in\mathfrak X^*(M)$。
\end{enumerate}
\end{proposition}

余向量场在流形上的几何解释不如向量场那样直观。在每个点 $p$ 处，余向量 $\omega_p\in T_p^*M$ 携带两条信息：它的核，即 $T_pM$ 中的一个线性超平面（余维 $1$ 线性子空间）；以及使 $\omega_p(v)=1$ 的向量 $v$ 集合，即一个平行于核的仿射超平面（\cref{figure:ch11-余向量场}）。（实际上，单是使 $\omega_p(v)=1$ 的集合就足够了，但把两个平行超平面都画出来有助于可视化。）任意其他向量 $v$ 的值 $\omega_p(v)$ 随后由线性插值或外插得到。

\begin{figure}[htbp]
\centering
\includegraphics[width=0.6\textwidth]{fig-11-1.png}
\caption{一个余向量场}\label{figure:ch11-余向量场}
\end{figure}

于是，可以把余向量场设想为在每个切空间中定义一对超平面，一条过原点、另一条平行于它，并随点连续变化。余向量场小的地方，其中一个超平面离核很远，最终在余向量场取零值的点处完全消失。
\end{document}
